% Source PDF pp. 31--36. The first continued sentence belongs to en-05.
In view of point a), to prove that $A=B$, it then suffices to show that one has:
\[
 \rk\mathfrak u-\rk\Centr_{\mathfrak u}(X')
 \leq\dim U-\dim\uCentr_U(X').
\]
Now since $U$ is smooth, one has $\dim U=\rk\mathfrak u$. On the other hand, one has (Exp. II 5.3.3):
\[
 \dim\uCentr_U(X')\leq\rk\Centr_{\mathfrak u}(X'),
\]
whence the result (note that one has in fact $\dim\uCentr_U(X')=\rk\Centr_{\mathfrak u}(X')$, which proves anew that $Z$ is smooth (5.7.1)).
% typo: the source's closing parenthesis after the parenthetical observation is missing; supplied.

\par\medskip\noindent\textbf{5.8. End of the proof of 5.1.1 i).}\ ---\par\label{XVII.5.8}
It remains for us to prove i) c) and i) d).

\noindent\emph{Proof of 5.1.1 i) d) ($U$ smooth, $H$ connected).} We shall need the following lemma:
\begin{lemma}{5.8.1}\label{XVII.5.8.1}
With the notation of 5.1.1, suppose that $U$ is smooth and $H$ radicial, and let $H_1$ be an algebraic subgroup of $H$ having a lifting $H'_1$ into $E$. Then $H$ has a lifting $H'$ into $E$ containing $H'_1$.
\end{lemma}
By induction on the height of $H/H_1$, one can assume that $H/H_1$ has height $1$. Let $C'=\uCentr_E(H'_1)$. I say that the canonical morphism $C'\to H$ is an epimorphism. To establish this point, we can assume that $k$ is algebraically closed; but then the extension $E$ is trivial (5.1.1 i a)); let $H''$ be a lifting of $H$ into $E$, and $H''_1$ the inverse image of $H_1$ in $H''$. The groups $H'_1$ and $H''_1$ are two liftings of $H_1$ into $E$, hence are conjugate by an element of $U(k)$, since $k$ is algebraically closed and $U$ smooth (5.1.1 ii b)). It is then clear that to prove the assertion about $C'$, it suffices to prove it for $C''=\uCentr_E(H''_1)$. But in this case, $C''$ contains $H''$, and the property is clear. Moreover, it follows from 5.7.1 that $C\cap U$ is smooth. It is then clear that we can replace $E$ by $C$, hence assume that $H'_1$ is central. But then, passing to the quotient by $H'_1$, we can assume $H_1=0$ and $H$ of height $1$.
% typo?: the source changes C' to C in the last part of this reduction; kept as printed.

Since $U$ is smooth, one has the exact sequence of Lie $p$-algebras (App. II 3.2):
\[
 0\to\Lie U\to\Lie E\to\Lie H\to0.\tag{*}
\]
In view of App. II 2.2, saying that $E$ is trivial is equivalent to saying that (*) is a trivial extension of Lie $p$-algebras. Suppose $H\ne0$ (hence $\Lie H\ne0$), and suppose that we have found a Lie $p$-subalgebra $\mathfrak h_1$ of $\mathfrak h=\Lie H$ that is nonzero and lifts to a Lie $p$-subalgebra $\mathfrak h'_1$ of $\Lie E$. By \loccit, there is a subgroup $H_1$ of $H$ such that $\Lie H_1=\mathfrak h_1$, and a lifting $H'_1$ of $H_1$ into $E$ such that $\Lie H'_1=\mathfrak h'_1$. Applying again the reduction described above, one reduces to the same problem, where $H$ has been replaced by $H/H_1$. Since $H$ has height $1$, $\Lie(H/H_1)=\Lie H/\Lie H_1$ (\loccit), hence $\rk\Lie(H/H_1)\leq\rk\Lie(H)-1$. In brief, proceeding by induction on the rank of $\Lie H$, one sees that it suffices, when $\mathfrak h\ne0$, to find a nonzero Lie subalgebra of $\Lie H$ that lifts into $\Lie E$. Now one has the following lemma:
\begin{lemma}{5.8.2}\label{XVII.5.8.2}
Let $k$ be a field of characteristic $p>0$, and $\varphi$ a surjective morphism of Lie $p$-$k$-algebras of finite rank $\mathfrak g\to\mathfrak h$. Then:
% Source PDF p. 32.
\begin{enumerate}[label=\textup{\roman*)}]
\item If $\mathfrak h_{\bar k}$ is reductive (4.2.2) and $\ne0$, there is a reductive Lie subalgebra $\mathfrak h_1$ of $\mathfrak g$ whose image in $\mathfrak h$ is nonzero.
\item If $k$ is perfect and if $u$ is a unipotent element of $\mathfrak h$ (i.e. there is $n$ such that $u^{(p^n)}=0$), then $u$ lifts to a unipotent element of $\mathfrak g$.
\end{enumerate}
\end{lemma}
Take for $X$ an element $\ne0$ of $\mathfrak h$ in case i) and $u$ in case ii), and let $X'$ be a lifting of $X$ in $\mathfrak g$. The Lie $p$-subalgebra of $\mathfrak g$ generated by $X'$ is an abelian Lie $p$-algebra (Exp. VII) $\mathfrak h'$.

\noindent\emph{Case i).} It is clear from the description given in 4.2.2 that the reductive part (\loccit) of $\mathfrak h'_{\bar k}$ is already defined over $k$; denote it by $\mathfrak r'$. I say that the image of $\mathfrak r'$ in $\mathfrak h$ is nonzero. To establish this point, we can assume that $k$ is algebraically closed, so that $\mathfrak h'=\mathfrak r'\oplus\mathfrak u'$ ($\mathfrak u'$ the unipotent part of $\mathfrak h'$). If the image of $\mathfrak r'$ in $\mathfrak h$ were zero, the image of $\mathfrak h'$ in $\mathfrak h$ would be unipotent, hence would be zero since $\mathfrak h$ is reductive (\cf 2.4 ii)), whereas by construction it contains $X$.

\noindent\emph{Case ii).} One proceeds in the same way, interchanging the roles of $\mathfrak r'$ and $\mathfrak u'$.

\noindent\emph{End of the proof of 5.8.1.} Assuming $H\ne0$, there is, by 5.8.2, a nonzero reductive Lie subalgebra $\mathfrak h'_1$ of $\Lie E$. Since $\Lie U$ is unipotent (4.3 i)), one necessarily has $(\Lie U)\cap\mathfrak h'_1=0$, so that $\mathfrak h'_1$ is a lifting of a Lie $p$-subalgebra of $\Lie H$.

\noindent\emph{End of the proof of 5.1.1 i d).} By 5.8.1, there is a family of algebraic subgroups $H'_n$ ($n\in\NN$) of $E$, such that $H'_{n+1}$ contains $H'_n$ and such that $H'_n$ lifts ${}_{F^n}H$. The decreasing sequence of subgroups $\uCentr_E(H'_n)$ is stationary; let $C$ be its stationary value. The center $Z$ of $C$ contains $H'_n$ for every $n$, hence the image of $Z$ in $H$ contains ${}_{F^n}(H)$ for every $n$, and consequently is an open subgroup of $H$ (Exp. VIIA \S4), hence is equal to $H$ since $H$ is connected. To prove that $E$ is trivial, one can therefore replace $E$ by $Z$, hence assume that $E$ is commutative. We shall then see in 7.2.1 that $E$ contains a maximal subgroup of multiplicative type $M$ whose formation commutes with extension of the base field. Since $E_{\bar k}$ is a trivial extension (5.1.1 i) a)) and $U$ is unipotent, it is clear that $M$ is the unique lifting of $H$ into $E$.

\noindent\emph{Proof of 5.1.1 i) c) ($U$ $k$-solvable).} Since $(H/H)^0(k)$ is of order prime to $p$, it is immediate by duality that there is an integer $n$ such that ${}_nH$ is an étale subgroup and such that the canonical morphism ${}_nH\to H/H^0$ is an epimorphism, so that $H/{}_nH$ is connected. By 5.2.3 d), there is a lifting $H'$ of ${}_nH$ into $E$. One shows, as at the beginning of the proof of 5.8.1, that $C=\uCentr_E(H')$ is a subgroup of $E$ such that $C\to H$ is an epimorphism and such that $C\cap U$ is smooth. Replacing $E$ by $C$ and passing to the quotient by $H'$, one is reduced to the case where $U$ is smooth and $H$ connected, that is, to case i) d).
% typo?: the source prints (H/H)^0(k), rather than (H/H^0)(k); kept as printed.

This completes the proof of 5.1.1.

\par\medskip\noindent\textbf{5.9. Counterexamples.}\ ---\par\label{XVII.5.9}
Let us first indicate a procedure for obtaining nontrivial extensions of a $k$-group of multiplicative type $H$ by a unipotent $k$-group $U$. Suppose that an action of $H$ on $U$ is given, and let $E$ be the semidirect product $E=U\cdot H$. Let on the other hand an element of $H^1(k,U)$ be represented by a $1$-cocycle $a$.
% Source PDF p. 33.
The group $U$ acts by inner automorphisms on $E$. The datum of $a$ therefore defines a $k$-form of $E$, denoted by $E_a$. Suppose moreover that $U$ is commutative; then $U$ acts trivially on $U$ and on the quotient $E/U=H$, so that $E_a$ is again an extension of $H$ by $U$. Suppose, for simplicity, that $H$ is an étale diagonalizable group, so that the Galois group $\mathscr G$ of $\bar k/k$ acts trivially on $H(k)$; the action of $\mathscr G$ on the points of $E_a(\bar k)$ is then given by the formula:
\[
 {}^g(u,h)=({}^gu+a(g)-{}^ha(g),h)
 \qquad(g\in\mathscr G,\ u\in U(\bar k),\ h\in H(\bar k)).
\]
If $h$ is a point of $H(k)$, and $X$ its inverse image in $E_a$, $X$ is therefore a torsor under $U$ defined by the class of the $1$-cocycle of $\mathscr G$ with values in $U$, $g\mto a(g)-{}^ha(g)$. It follows that if there is a point $h$ of $H(k)$ such that the preceding $1$-cocycle is nontrivial, the extension $E_a$ is nontrivial. We shall apply this construction in two particular cases:

\noindent a) \emph{Nontrivial extension of an étale diagonalizable group $H$ by $\ZZ/p\ZZ$.}

Take for $H$ the group $(\ZZ/p\ZZ)^\times$, which one makes act by multiplication on $U=\ZZ/p\ZZ$. Let on the other hand $k$ be a field of characteristic $p$, $K$ an extension of $k$ with Galois group $\mathscr G$ isomorphic to $\ZZ/p\ZZ$, and $a$ a nonzero element of $H^1(\mathscr G,U)=\Hom(\mathscr G,U)$. The group $E_a$ then answers the question.

\noindent b) \emph{Example of a nontrivial extension of an étale diagonalizable group $H$ by a smooth connected unipotent group $U$.}\par\nopagebreak[4]

Take for the field $k$ a nonperfect field such that there exists a $k$-form $U$ of $\Ga$ such that $H^1(k,U)\ne0$. For example (\cf J.-P. Serre, \emph{Cohomologie Galoisienne}), one can take for $k$ the fraction field of a discrete valuation ring of equal characteristic $p>0$, and for $U$ the algebraic subgroup of $\Ga\times_k\Ga$ with equation $X^p+X+tY^p=0$, where $t$ denotes a uniformizer of $k$. Indeed, assuming for simplicity that $k$ contains the $(p-1)$th roots of unity, one has an exact sequence of algebraic groups:
\[
 \begin{array}{ccccc}
 0\to\ZZ/p\ZZ&\to&U&\to&\Ga\to0\\
 &&(x,y)&\mto&y,
 \end{array}
\]
hence an exact cohomology sequence:
\[
 \Ga(k)\xrightarrow{d}H^1(k,\ZZ/p\ZZ)\to H^1(k,U)\to0,
\]
where $d$ associates to $x$ the homogeneous space under $\ZZ/p\ZZ$ with equation
\[
 X^p+X+tx^p=0.
\]
Moreover, one knows that $H^1(k,\ZZ/p\ZZ)$ is isomorphic to $k/\mathfrak P(k)$ (where $\mathfrak P(x)=x^p+x$), hence $H^1(k,U)$ is isomorphic to $k/(\mathfrak P(k)+tk^p)$. Suppose moreover $p\ne2$; it is then clear that $t^2$ is an element of $k$ that does not belong to $\mathfrak P(k)+tk^p$, hence $H^1(k,U)\ne0$.
% typo?: the source uses x^p+x and t^2, with the stated root-of-unity hypothesis; kept as printed.
On the other hand $\bmu_{p-1}$ acts on $U$ by the formula:
\[
 (h,x,y)\mto(hx,hy).
\]
Denote by $\mathscr G$ the Galois group of the extension $K$ defined by the equation
\[
 X^p+X+t^2=0,
\]
% Source PDF p. 34.
and let $a\in H^1(\mathscr G,U)$ be the nonzero element described above. One immediately checks that $E_a$ is then a nontrivial extension of $\bmu_{p-1}$ by $U$.

\noindent c) \emph{Nontrivial extension of $\Gm$ by $\balpha_p$.}

By 5.1.1 i) b), such an extension can exist only over a nonperfect field $k$. Let therefore $k$ be a nonperfect field, $G$ the semidirect product of $U=\Ga$ by $H=\Gm$, with $\Gm$ acting on $\Ga$ by homotheties. Since $U$ is invariant in $G$, ${}_FU$ is also. Let $G'=G/{}_FU$. The group $G'$ is isomorphic to $\Ga\cdot\Gm$, where $\Gm$ acts on $\Ga$ by the formula:
\[
 (h,u)\mto h^pu.
\]
The functor $\mathscr T_G$ (resp. $\mathscr T_{G''}$) of subtori of $G$ (resp. $G''$) (\cf Exp. XV) is isomorphic to $\Ga$, and the morphism $\mathscr T_G\to\mathscr T_{G''}$ deduced from the morphism $G\to G''$ is identified with the morphism $u\mto u^p$. It follows that if $T''$ is a subtorus of $G''$ corresponding to a point $x$ of $k\simeq\Ga(k)$ such that $x^{1/p}$ is not in $k$, the inverse image $E$ of $T''$ in $G$ will be an extension of a torus $T''$ by ${}_FU=\balpha_p$, will have no maximal tori defined over $k$, hence will be nontrivial. One finds for $E$ the subgroup of $G=\Spec[U,T,T^{-1}]$ with equation $U^p=x-xT^p$.
% typo?: the source defines G' and then uses G'' throughout the torus-functor discussion; kept as printed.
% typo?: the source prints Spec[U,T,T^{-1}] without a coefficient field; kept as printed.

\begin{remark}{5.9.1}\label{XVII.5.9.1}
This last example shows that a nonsmooth algebraic group defined over a nonperfect field does not necessarily have maximal tori, and thus answers the question asked in Exp. XIV 1.5 b).
\end{remark}
\noindent d) \emph{Let us now give an example of a trivial extension $E$ of a group of multiplicative type $H$ by a unipotent group $U$, and of two liftings $H'$ and $H''$ of $H$ that are not conjugate by an element of $U(k)$.}

Take for $E$ the semidirect product of $U=\Ga$ by $\bmu_p=\Spec k[T]/(T^p-1)$, the action of $\bmu_p$ on $\Ga=\Spec k[U]$ being defined by the comorphism:
\[
 U\mto(U+U^p)T-U^p.
\]
The centralizer of $\bmu_p$ in $U$ is then the étale group $Z$ with equation $U+U^p=0$. It follows that if $k$ is not algebraically closed, the canonical map $U(k)\to(U/Z)(k)$ is not in general surjective, hence, in view of 5.1.1 ii) b), two liftings of $\bmu_p$ into $E$ are not necessarily conjugate by an element of $U(k)$.

Here is another example, with $k$ algebraically closed of characteristic $p>0$. Let $G$ be the radicial group that is the semidirect product of $\bmu_p$ by $\balpha_p$, where $\bmu_p$ acts on $\balpha_p$ by ``homotheties''. One then has an exact sequence of algebraic groups with operator group $\bmu_p$:
\[
 \begin{array}{ccccccc}
 0&\to&\balpha_p&\to&\Ga&\to&\Ga\to0\\
 &&&&x&\mto&x^p,
 \end{array}
\]
where $\bmu_p$ acts by homotheties on the first term $\Ga$, and trivially on the second. The exact cohomology sequence (App. I, prop. 11) gives here the exact sequence:
\[
 0\to\Ga(k)\to H^1(\bmu_p,\balpha_p)\to H^1(\bmu_p,\Ga).
\]
Since the last term is zero (I 5.3.3), one sees that $H^1(\bmu_p,\balpha_p)$ is nonzero, hence two liftings of $\bmu_p$ into $G$ are not necessarily conjugate.

% Source PDF p. 35.
\sgasection{6}{Extension of a unipotent group by a group of multiplicative type}
\par\medskip\noindent\textbf{6.1. Statement of the theorem.}\ ---\par\label{XVII.6.1}
\begin{theorem}{6.1.1}\label{XVII.6.1.1}
Let $k$ be a field, $U$ a unipotent algebraic $k$-group, $H$ a $k$-group of multiplicative type, and $E$ an algebraic $k$-group that is an extension of $U$ by $H$, so that one has the exact sequence
\[
 1\to H\to E\to U\to1.
\]
Then the extension $E$ is trivial and there is a unique lifting of $U$ into $E$ in each of the following cases:
\begin{enumerate}[label=\textup{\Alph*)}]
\item The group $U$ is smooth and one of the following conditions is satisfied:
\begin{enumerate}[label=\textup{\roman*)}]
\item $U$ is connected and the canonical morphism $E\to U$ has a section.
\item $U$ has a composition series with successive quotients isomorphic to $\Ga$.
\item $H$ is étale.
\item $k$ is perfect.
\end{enumerate}
\item $U=\balpha_p$ and $k$ is perfect.
\item $E$ is commutative and $k$ is perfect.
\end{enumerate}
\end{theorem}

\par\medskip\noindent\textbf{6.2. Proof of 6.1.1 A).}\ ---\par\label{XVII.6.2}
Let us first establish three lemmas.
\begin{lemma}{6.2.1}\label{XVII.6.2.1}
Let $S$ be a prescheme, $E$ a group $S$-prescheme that is an extension of a group $S$-prescheme $U$ with connected fibers by an $S$-group of multiplicative type and of finite type $H$ (i.e. $U$ is the quotient of $E$ by $H$ for the fpqc topology). Then $E$ is a central extension.
\end{lemma}
Indeed, since $H$ is commutative, the group $U$ acts by inner automorphisms on $H$, through a group $S$-morphism
\[
 u:U\to\underline{\Aut}_{S\text{-gr}}(H).
\]
The functor $\underline{\Aut}_{S\text{-gr}}(H)$ is representable by an étale $S$-scheme (Exp. X 5.10), and consequently the identity section is an immersion that is both open and closed. Since $U$ has connected fibers, one deduces that $u$ is the unit morphism.

\begin{lemma}{6.2.2}\label{XVII.6.2.2}
With the notation of 6.1.1, if $E$ is trivial, there is a unique lifting of $U$ into $E$.
\end{lemma}
Let therefore $U'$ and $U''$ be two liftings of $U$ into $E$. To show that $U'=U''$, we can assume that $k$ is algebraically closed, and it suffices to show that $H^1(U,H)=0$. If $U$ is connected, $U$ centralizes $H$ (6.2.1), hence
\[
 H^1(U,H)=\Hom_{k\text{-gr}}(U,H)=0
\]
by 2.4 ii). In the general case, denote by $U'_1$ the unique lifting into $E$ of the connected component $U^0$ of $U$, and let $N=\uNorm_E(U'_1)$. If $g\in E(k)$, $\operatorname{int}(g)U'_1$ is a lifting of $U^0$, hence is equal to $U'_1$, and consequently $N(k)\supset E(k)$. Moreover, $N$ contains $H$ (6.2.1) and $U'_1$, whence immediately the fact that $N=E$. Passing to the quotient by $U'_1$, one is reduced to the case where $U$ is étale. In this case, since $k$ is algebraically closed, $H^1(U,H)$ is identified with the ordinary cohomology group $H^1(U(k),H(k))$\nde{See, for example, Proposition III.6.4.2 of the book by M. Demazure and P. Gabriel, \emph{Groupes algébriques I}, Masson \& North-Holland (1970).} and consequently is zero, since $U(k)$ is a finite $p$-group and $H(k)$ is uniquely $p$-divisible.
% Source PDF p. 36.
% typo: editor note 10 prints "(1970."; supplied its closing parenthesis.

\begin{corollary}{6.2.3}\label{XVII.6.2.3}
To prove that $E$ is trivial, it suffices to show that $E$ becomes trivial after finite separable extension of the base field; in particular, one can assume that $H$ is diagonalizable.
\end{corollary}
\begin{lemma}{6.2.4}\label{XVII.6.2.4}
For every central extension $E$ of $U$ by a diagonalizable $k$-group $H$ to be trivial, it suffices that every central extension of $U$ by $\Gm$ be trivial.
\end{lemma}
Indeed, by induction on $r$, one first notes that the hypothesis made implies that $E$ is trivial if $H=\Gm^r$. In the general case, $H$ embeds into $\Gm^r$ for a suitable integer $r$ (this is immediate by duality); let $H''=\Gm^r/H$. One obtains the exact sequence (App. I 2.1):
\[
 Z^1(U,H'')\to\Ext_{\mathrm{alg}}(U,H)\to\Ext_{\mathrm{alg}}(U,\Gm^r)=0
\]
(where $U$ acts trivially on $H$, $H''$, $\Gm^r$). But $Z^1(U,H'')=\Hom_{k\text{-gr}}(U,H'')=0$ (2.4 ii)), hence $\Ext_{\mathrm{alg}}(U,H)=0$.

\noindent\emph{Proof of 6.1.1 A) i).} Since $E\to U$ has a section, $E$ defines an element $\bar f$ of $H^2(U,H)$ (App. I 3.1). We must show that $\bar f$ is zero, and for this it suffices to show that a $2$-cocycle $f:U\times_k U\to H$ is a constant morphism, which will follow from the following lemma:
\begin{lemma}{6.2.5}\label{XVII.6.2.5}
Let $U$ be a smooth connected unipotent algebraic $k$-group, and $H$ a $k$-group of multiplicative type; then every $k$-morphism (of preschemes):
\[
 f:U\to H
\]
is constant.
\end{lemma}
To prove this lemma, we can assume that $k$ is algebraically closed. We proceed by increasing induction on $\dim U$. If $\dim U>0$, $U$ has a composition series (\cf 3.9):
\[
 1\to U'\to U\xrightarrow{\pi}U''\to1
\]
with $U'\simeq\Ga$ and $\dim U''<\dim U$. It suffices to show that $f$ factors through $U''$. Since the graph of the equivalence relation defined by $\pi$ is smooth over $k$, hence reduced, it suffices to show that if $x,y\in U(k)$ have the same image $z$ in $U''(k)$, then $f(x)=f(y)$. Now $\pi^{-1}(z)$ is isomorphic to $\Ga$, hence the restriction of $f$ to $\pi^{-1}(z)$ factors through a reduced irreducible component of $H$, hence through a $k$-scheme isomorphic to $(\Gm)^r$. It then suffices to note that every morphism from $\Ga$ to $(\Gm)^r$ is constant, since every invertible regular function on $\Ga$ is constant.

\noindent\emph{Proof of 6.1.1 A) ii).} By 6.2.1, 6.2.3 and 6.2.4, we can assume that $H=\Gm$. By hypothesis, $U$ has a composition series $U\supset U_1\supset\cdots$ such that $U_i/U_{i+1}$ is isomorphic to $\Ga$. Let $E_1$ be the inverse image of $U_1$ in $E$. By induction on $\dim U$, one can assume that the extension $E_1$ is trivial; let $U'_1$ be the unique lifting of $U_1$.
% The last sentence begins on p. 36 and finishes on p. 37; en-07 starts "Proceeding as in the proof of 6.2.4". No open environment.
